% Part: lambda-calculus
% Chapter: introduction
% Section: composition

\documentclass[../../../include/open-logic-section]{subfiles}

\begin{document}

\olfileid[pt]{lam}{int}{com}
\olsection{As funções \usetoken{S}{lambda definable} são fechadas sob composição}

\begin{lem}
As funções !!{lambda definable} são fechadas sob composição.
\end{lem}

\begin{proof}
Suponha que $f$ seja definida por composição a partir de $h$, $g_0,$
\dots,~$g_{k-1}$. Supondo que $h$, $g_0$, \dots,~$g_{k-1}$ sejam
!!{lambda defined} por $H$, $G_0$, \dots,~$G_{k-1}$, respectivamente,
precisamos encontrar um termo~$F$ que !!{lambda define}s~$f$. Mas podemos
simplesmente definir~$F$ por
\[
F(x_0, \dots, x_{l-1}) = H(G_0(x_0, \dots, x_{l-1}),
\dots, G_{k-1}(x_0, \dots, x_{l-1})).
\]
Em outras palavras, a linguagem do cálculo lambda é bem adequada para
representar a composição.
\end{proof}

\end{document}
